% 18 November 2026. Prepared for Edward Baker with OpenAI Codex assistance.
% Actual recorded model and effort unavailable. B10 formulation, E31 integration.
\section{Recursive and continuous two-sided refinement}
\label{sec:two-sided}

The preceding sections suggest a more informative use of forward and
backward propagation. A forward source calculation may constrain the
environment of a small gap; a backward comparison may then reduce the
admissible earlier configurations; the reduced configuration set may
sharpen the next forward calculation. The object being refined is an
\emph{enclosure of one source-consistent trajectory}, not the physical
trajectory itself. We develop the formulation proposed in
\cite{BakerCM35}. Its elementary soundness and residual identities are
proved below. Strict improvement for the canonical source remains an
additional quantitative question.

There are two distinct parameters: physical Newman time $t$, and the
iteration number $k$ of an enclosure calculation. Introducing $k$ does
not reverse the physical heat equation, generate another divisor or
supply a finite moment closure. The complete source, the label/count
record and the errors must be shared by the two directions. In
particular, Theorem~\ref{rf:horizon} continues to screen repetition
of the unchanged initial-interaction consumer within its stated class.

\subsection{Common-source state sets and sound refinement}

Fix finite times $t_0<\cdots<t_m$ inside independently cleared counted
contours. Write $\omega$ for a common complete-source record, including
normalization, shared source-error parameters and count/label information.
Let $z_i$ consist of regular observables at $t_i$. For a separated pair
these may include $c,D,M$ and source traces; their admissible domain must
allow both signs of $D$ if positivity is the desired conclusion. For an
overlapping block, choose a valid common chart or larger-cluster symmetric
coordinates when a neighboring edge chart fails. An unexcluded chart or
boundary exit remains a possible failure; it cannot be discarded in the
definition of the admissible set.

Let $E_i^0$ be an outer set of augmented states $(\omega,z_i)$ containing
the true $(\omega_*,z_i^*)$. A validated transition relation
$\mathcal R_i$ contains the true transition and relates only states with
the same source record:
\begin{equation}
 ((\omega,z_i),(\omega,z_{i+1}))\in\mathcal R_i.
 \label{ts:transition}
\end{equation}
This relation may be an outer approximation. It is not asserted to
characterize exactly which finite moment data arise from a canonical
divisor. Define its two projections by
\begin{align}
 \operatorname{Post}_i(A)
 &=\{v:\text{there is }u\in A\text{ with }(u,v)\in\mathcal R_i\},
 \notag\\
 \operatorname{Pre}_i(B)
 &=\{u:\text{there is }v\in B\text{ with }(u,v)\in\mathcal R_i\}.
 \label{ts:projections}
\end{align}
Any new observation set $O_i^k$ must independently contain the true
augmented state. Take it to be the whole state space when no new source
observation is available. One simultaneous sweep is
\begin{equation}
 E_i^{k+1}=E_i^k\cap O_i^k
       \cap\operatorname{Post}_{i-1}(E_{i-1}^k)
       \cap\operatorname{Pre}_{i}(E_{i+1}^k),
 \label{ts:sweep}
\end{equation}
where nonexistent endpoint factors are omitted.

\begin{proposition}[Sound two-sided refinement]\label{ts:sound}
For every finite $k$, the sets in \eqref{ts:sweep} are nested and contain
the true augmented states. Define the full skeleton set
\begin{equation}
 \mathcal P^k=\{(\omega,z_0,\ldots,z_m):
       (\omega,z_i)\in E_i^k\text{ for every }i,
       \ ((\omega,z_i),(\omega,z_{i+1}))\in\mathcal R_i
                         \text{ for every }i<m\}.
 \label{ts:skeleton}
\end{equation}
Then $\mathcal P^{k+1}\subseteq\mathcal P^k$, and every $\mathcal P^k$
contains the true skeleton. For a fixed real statistic $Q$, wherever
the bounds are finite,
\begin{equation}
 \inf_{\mathcal P^{k+1}}Q\ge\inf_{\mathcal P^k}Q,
 \qquad
 \sup_{\mathcal P^{k+1}}Q\le\sup_{\mathcal P^k}Q.
 \label{ts:monotone-bounds}
\end{equation}
More precisely, for the fixed relations in \eqref{ts:transition},
\begin{equation}
 \mathcal P^{k+1}=\{p\in\mathcal P^k:
           (\omega,z_i)\in O_i^k\text{ for every }i\}.
 \label{ts:skeleton-invariance}
\end{equation}
Thus if all observation sets are universal, the full skeleton set and
its exact statistic bounds are unchanged by a sweep.
\end{proposition}

\begin{proof}
Assume the true states belong to all $E_i^k$. Their true neighboring
transitions witness membership in both projections in
\eqref{ts:sweep}, and soundness of $O_i^k$ supplies the remaining new
factor. The intersection therefore retains the true state and is a
subset of $E_i^k$. Induction starts at $k=0$. Substitution into
\eqref{ts:skeleton} proves its corresponding containment and
nonemptiness. Infima increase and suprema decrease when a nonempty
feasible set is restricted, proving \eqref{ts:monotone-bounds}.
For \eqref{ts:skeleton-invariance}, one inclusion follows from the
definition. Conversely, each skeleton in $\mathcal P^k$ supplies its
own neighboring witnesses for every projection in \eqref{ts:sweep}.
If it also satisfies the new observations, all its nodes survive.
\end{proof}

The proposition proves nonworsening enclosures, not strict contraction,
finite termination, a convergence rate or convergence to the exact
feasible set. An implementable projection may use an outward hull of
the exact image; intersecting that hull with the previous enclosure
remains sound if the true image is retained. Projection sweeps alone can
tighten marginal outer estimates by removing states without compatible
neighbors, but cannot improve the exact joint bounds over the already
retained skeleton. New observations or soundly refined transition
relations are needed to shrink that full feasible set. Outward hulls can
lose most of the marginal improvement. Similarly, discarding $\omega$ and choosing
unrelated forcing values on successive cells loses the common-source
constraint. Correlations across nodes and times belong in
\eqref{ts:skeleton}, even when the final reported statistic is scalar.

A statistic of a whole interval needs full admissible cell paths or a
paid between-node interpolation bound. For example, a lower bound for
$\min_{j,t}D_j(t)$ is not obtained from positive values at finitely many
times. A failure-count statistic may include chart exits, block-boundary
events and multiplicities; an upper bound below one then excludes those
events only after the population covers every eligible failure. Neither
random labeling nor a small average supplies this coverage.

\subsection{What reversal alone cannot add}

\begin{proposition}[Exact-flow consistency]\label{ts:inverse}
For a prescribed comparison flow $\Phi_{b,a}$ on a domain where the
forward and reverse solutions are unique,
\begin{equation}
 \Phi_{a,b}\bigl(\Phi_{b,a}(A)\bigr)=A
 \label{ts:inverse-identity}
\end{equation}
for every initial set whose trajectories stay in that domain.
\end{proposition}

\begin{proof}
Uniqueness identifies the reverse of each forward trajectory with its
original trajectory. Thus the two maps compose to the identity on the
stated initial set.
\end{proof}

This is a consistency test for an exact prescribed flow, not an
impossibility theorem for source acquisition. A terminal set obtained
only as the exact forward image of the initial set gives no extra
constraint upon exact reversal. In contrast, interval propagation may
have discarded dependencies which a second sweep recovers, or a later
source calculation may supply an additional checked observation.
Both are legitimate reasons for shrinking an outer set. They must be
compared with direct joint propagation using the same information.
No favorable stability of the inverse is implied by
\eqref{ts:inverse-identity}.

\subsection{An exact continuous forward--adjoint pairing}

For a regular counted pair, retain physical forward time and write
\begin{equation}
 y=\binom D M,\qquad
 y'=G_{\mathcal F}(t,y)
   =\binom{2-4DM}{2M^2-\mathcal F(t)}.
 \label{ts:pair}
\end{equation}
The function $\mathcal F(t)$ denotes the complete forcing along the
\emph{same actual source-consistent path}. Equation~\eqref{ts:pair}
does not say that two moments determine that forcing. Its dependence on
the source, center and exterior remains in $\omega$ and must be acquired
by the contour or hierarchy arguments already developed.

Choose a $C^1$ reference path $\widehat y$ and continuous reference
forcing $\widehat{\mathcal F}$. Put
\begin{equation}
 e=y-\widehat y,\quad
 \delta\mathcal F=\mathcal F-\widehat{\mathcal F},\quad
 r=\widehat y'-G_{\widehat{\mathcal F}}(t,\widehat y),\quad
 \widehat A=
 \begin{pmatrix}-4\widehat M&-4\widehat D\\0&4\widehat M\end{pmatrix}.
 \label{ts:reference}
\end{equation}
The complete error equation is
\begin{equation}
 e'=\widehat A e+
       \binom{-4e_De_M}{2e_M^2-\delta\mathcal F}-r.
 \label{ts:error}
\end{equation}
Indeed, expanding $(\widehat D+e_D)(\widehat M+e_M)$ and
$(\widehat M+e_M)^2$ gives every term displayed in
\eqref{ts:error}. The quadratic remainder has not been dropped.

\begin{proposition}[Two-sided adjoint residual identity]\label{ts:adjoint}
On a compact regular interval $[a,b]$, fix $\tau\in[a,b]$ and a vector
$\ell\in\R^2$. Let
\begin{equation}
 p'=-\widehat A^{\mathsf T}p,\qquad p(\tau)=\ell,
 \qquad
 \Psi=-4p_De_De_M+2p_Me_M^2-p_M\delta\mathcal F-p\cdot r.
 \label{ts:adjoint-def}
\end{equation}
Then
\begin{align}
 \ell\cdot e(\tau)
 &=p(a)\cdot e(a)+\int_a^\tau\Psi(t)\dd t\notag\\
 &=p(b)\cdot e(b)-\int_\tau^b\Psi(t)\dd t.
 \label{ts:two-sided}
\end{align}
The identity has no denominator $D$ and is valid through a target merger
whenever the complete counted-pair tube and regular quantities remain
defined.
\end{proposition}

\begin{proof}
The continuous coefficient matrix gives a unique linear adjoint solution
on the compact interval in both time directions. Apply the product rule
to $p\cdot e$, substitute \eqref{ts:error} and
\eqref{ts:adjoint-def}, and cancel the two linear terms. The result is
$(p\cdot e)'=\Psi$. Integration separately from $a$ to $\tau$ and from
$\tau$ to $b$, with $p(\tau)=\ell$, proves
\eqref{ts:two-sided}. All expressions remain polynomial in the pair
errors at $D=0$; contour separation, not division by $D$, is the domain
condition here.
\end{proof}

Taking $\ell=(1,0)$ reads out the discriminant error. If paid enclosures
of the two right sides are intervals $I_L(\tau)$ and $I_R(\tau)$, then
\begin{equation}
 D(\tau)\in\widehat D(\tau)+
                  \bigl(I_L(\tau)\cap I_R(\tau)\bigr).
 \label{ts:intersection}
\end{equation}
This supplies a continuous counterpart to the discrete sweep: a sound
smaller state enclosure can be fed into the source residual and the
quadratic terms, and the resulting two readouts can be intersected again.
The intervals enclose two expressions for the same error. They are not
independent measurements. Disjoint intervals mean that the asserted
enclosures or their source assumptions are inconsistent; they do not
prove positivity of an excluded canonical path.

\subsection{Error payment, statistics and strict improvement}

For transparent sufficient error payment, suppose
$|e_D|\le E_D$, $|e_M|\le E_M$,
$|\delta\mathcal F|\le E_{\mathcal F}$ and $|r_i|\le R_i$.
Then
\begin{equation}
 |\Psi|\le
 4|p_D|E_DE_M+2|p_M|E_M^2+|p_M|E_{\mathcal F}
                    +|p_D|R_D+|p_M|R_M.
 \label{ts:absolute}
\end{equation}
This follows term by term from \eqref{ts:adjoint-def}. It is not a
claim that separate absolute allowances are efficient. The signed
correlation of the forcing error, reference defect and quadratic terms
may be essential. The complete source contribution to
$E_{\mathcal F}$ includes uncertain trial coordinates, with the kernel
derivatives of Section~\ref{sec:retained-forcing}, together with source
tails, contour/quadrature, arithmetic and whole-time remainders. Endpoint
and adjoint-solution errors must also be included when those objects are
computed rather than exact. Large adjoint coefficients can erase any
apparent improvement.

Fixed finite weighted sums of \eqref{ts:two-sided} give the corresponding
identity for linear empirical error statistics, with all labels and the
same source retained. They do not give a sharp bad-gap probability by
themselves. A bad-event majorant needs its own derivative/resolution
payment and complete population coverage; an integer zero-exception
argument needs a bound below the appropriate integer threshold. An
average over time can miss an isolated event. Whole-interval use of
\eqref{ts:intersection} therefore requires a uniform bound for every
$\tau$, paid interpolation on time cells, or a first-exit argument whose
other exits are independently excluded.

The two-sided construction distinguishes a possible route from an
acquired gain. It preserves the physical forcing discarded by the
initial-$M$ horizon and can preserve correlations lost by an unsigned
force norm. It has not yet supplied a strict canonical intersection
width, a better whole-time gap margin, a convergence rate or a lower
acquisition cost. In particular it does not acquire the forcing floor
used in Appendix~\ref{app:compact-example} merely by repeating a sweep.

Before a quantitative execution, one must fix a common finite pair or
block, physical interval, independently cleared source/count tube,
regular charts, endpoint observations, reference path and complete
signed residual budget. A proposed improvement must be strict and
measured against direct regular $D/M$ or source-discriminant evolution
with precisely the same information, including the elapsed forward
excursion. Finite work and refinement limits and a separate check belong
to that contract. The present section supplies the analytical interface,
not its missing source estimate.

Quadratic and cosine-based heat controls satisfy the local identities on
their valid charts. A sound enclosure whose input admits their collision
must retain that event. The canonical application uses the separate
source-equality, order/count, physical-readout and integer-multiplicity
premises summarized in Section~\ref{sec:discussion}; none is an automatic
favorable residual bound. Thus the recursive and continuous descriptions
are mathematically specified without being promoted to a proved recursive
improvement of Newman bounds.
